\section{Week 10: Bayesian Modeling II: Posterior Predictive Checks and
Bayes Factors}

\begin{quote}
\textbf{Chapter 4: Bayesian modeling} (with Chapter 3: evaluation) ·
Lecture 1 h · Lab 2 h \textbf{Running data:} the Week 2 forgetting data
(exponential vs.~power), and a weak-evidence SDT participant (is
\(d' = 0\)?).
\end{quote}

\begin{center}\rule{0.5\linewidth}{0.5pt}\end{center}

\subsection{Key terms}

Important terms of this lesson, with their Persian equivalents (repeated
even if they appeared in earlier weeks).

\begin{longtable}[]{@{}
  >{\raggedright\arraybackslash}p{(\linewidth - 4\tabcolsep) * \real{0.3333}}
  >{\raggedright\arraybackslash}p{(\linewidth - 4\tabcolsep) * \real{0.3333}}
  >{\raggedright\arraybackslash}p{(\linewidth - 4\tabcolsep) * \real{0.3333}}@{}}
\caption{\centering Key terms}\tabularnewline
\toprule\noalign{}
\begin{minipage}[b]{\linewidth}\raggedright
English
\end{minipage} & \begin{minipage}[b]{\linewidth}\raggedright
\RL{\CMfa فارسی}
\end{minipage} & \begin{minipage}[b]{\linewidth}\raggedright
\RL{\CMfa توضیح کوتاه}
\end{minipage} \\
\midrule\noalign{}
\endfirsthead
\toprule\noalign{}
\begin{minipage}[b]{\linewidth}\raggedright
English
\end{minipage} & \begin{minipage}[b]{\linewidth}\raggedright
\RL{\CMfa فارسی}
\end{minipage} & \begin{minipage}[b]{\linewidth}\raggedright
\RL{\CMfa توضیح کوتاه}
\end{minipage} \\
\midrule\noalign{}
\endhead
\bottomrule\noalign{}
\endlastfoot
posterior & \RL{\CMfa توزیع پسین} (posterior) & \RL{\CMfa باور درباره‌ی پارامتر بعد
از دیدن داده} \\
forgetting curve & \RL{\CMfa منحنی فراموشی} & \RL{\CMfa کاهش یادآوری با گذشت
زمان} \\
exponential / power function & \RL{\CMfa تابع نمایی} / \RL{\CMfa توانی} & \RL{\CMfa دو شکل
رایج برای منحنی فراموشی و یادگیری} \\
evidence & \RL{\CMfa شواهد} & \RL{\CMfa اطلاعاتی که برای تصمیم جمع می‌شود} \\
sensitivity (d′) & \RL{\CMfa حساسیت} & \RL{\CMfa فاصله‌ی توزیع علامت و نویز} \\
posterior predictive check (PPC) & \RL{\CMfa بررسی پیش‌بینی پسین} & \RL{\CMfa مقایسه‌ی
داده‌ی شبیه‌سازی‌شده از پسین با داده‌ی واقعی} \\
marginal likelihood / evidence & \RL{\CMfa درست‌نمایی حاشیه‌ای} & \RL{\CMfa احتمال
داده تحت مدل، میانگین روی پیشین} \\
likelihood & \RL{\CMfa درست‌نمایی} & \RL{\CMfa احتمال داده به‌عنوان تابعی از
پارامترها} \\
Bayes factor (BF) & \RL{\CMfa عامل بیز} (Bayes factor) & \RL{\CMfa نسبت درست‌نمایی‌های
حاشیه‌ای دو مدل} \\
model flexibility / complexity & \RL{\CMfa انعطاف‌پذیری} / \RL{\CMfa پیچیدگی مدل} &
\RL{\CMfa دامنه‌ی الگوهایی که مدل می‌تواند تولید کند} \\
grid approximation & \RL{\CMfa تقریب شبکه‌ای} (grid) & \RL{\CMfa محاسبه‌ی پسین روی
شبکه‌ای از مقادیر} \\
Savage--Dickey ratio & \RL{\CMfa نسبت} Savage--Dickey & BF \RL{\CMfa از نسبت چگالی
پسین به پیشین در مقدار صفر} \\
AIC / BIC & \RL{\CMfa معیار اطلاعاتی آکائیکه} / \RL{\CMfa بیزی} & \RL{\CMfa برازش به‌علاوه‌ی
جریمه برای تعداد پارامتر} \\
prior & \RL{\CMfa توزیع پیشین} (prior) & \RL{\CMfa باور درباره‌ی پارامتر پیش از دیدن
داده} \\
estimate / estimator & \RL{\CMfa برآورد} / \RL{\CMfa برآوردگر} & \RL{\CMfa مقدار به‌دست‌آمده
برای پارامتر} / \RL{\CMfa روش به دست آوردن آن} \\
posterior odds & \RL{\CMfa بخت پسین} (odds) & \RL{\CMfa نسبت احتمال پسین دو مدل} \\
fit / fitting & \RL{\CMfa برازش} & \RL{\CMfa تنظیم پارامترها تا مدل با داده جور
شود} \\
parameter & \RL{\CMfa پارامتر} & \RL{\CMfa کمیت قابل تنظیم مدل که بین افراد یا شرایط
فرق می‌کند} \\
simulation & \RL{\CMfa شبیه‌سازی} & \RL{\CMfa تولید داده‌ی مصنوعی با اجرای مدل} \\
noise & \RL{\CMfa نویز} & \RL{\CMfa تغییرات تصادفی در شواهد یا در داده} \\
adequacy check & \RL{\CMfa بررسی کفایت مدل} & \RL{\CMfa آیا مدل ویژگی‌های اصلی داده
را بازتولید می‌کند} \\
nested models & \RL{\CMfa مدل‌های تودرتو} (nested) & \RL{\CMfa یک مدل حالت خاص مدل
دیگر است} \\
trial & \RL{\CMfa کوشش} (trial) & \RL{\CMfa یک بار انجام تکلیف توسط شرکت‌کننده} \\
accuracy & \RL{\CMfa دقت} & \RL{\CMfa نسبت پاسخ‌های درست} \\
\end{longtable}

\begin{center}\rule{0.5\linewidth}{0.5pt}\end{center}

\subsection{Learning outcomes}

By the end of this week, students will be able to:

\begin{enumerate}
\def\labelenumi{\arabic{enumi}.}
\tightlist
\item
  \textbf{Generate} posterior predictive distributions and
  \textbf{conduct} graphical and numerical posterior predictive checks
  (PPCs).
\item
  \textbf{Define} the marginal likelihood and the Bayes factor, and
  \textbf{explain} why the marginal likelihood penalizes model
  flexibility automatically.
\item
  \textbf{Compute} Bayes factors by grid integration and with the
  Savage--Dickey density ratio, and \textbf{compare} them with the BIC
  approximation.
\item
  \textbf{Evaluate} the sensitivity of Bayes factors to the prior, and
  \textbf{contrast} it with the robustness of posterior estimates.
\item
  \textbf{Interpret} Bayes factors on a graded evidence scale and
  \textbf{combine} them with prior odds into posterior model
  probabilities.
\end{enumerate}

\begin{center}\rule{0.5\linewidth}{0.5pt}\end{center}

\subsection{Motivation \& big picture}

Last week gave us posteriors \emph{within} a model. Two questions
remain.

\textbf{1. Is the model any good at all?} A posterior can be narrow and
confident even when the model is badly wrong: the exponential model will
happily give a precise posterior for \(\lambda\) for the Week 2 data.
The Bayesian answer is to ask the fitted model to \textbf{generate new
data} and check whether they look like the real data. This is a
\emph{posterior predictive check}, the Bayesian counterpart of residual
plots and \(G^2\).

\textbf{2. Which model do the data support more?} In Week 7 we compared
models with AIC and BIC, which penalize parameter counts. Bayesian
inference has a built-in answer: compute how well each model predicted
the data \emph{before} it was fitted, averaging over all parameter
values its prior allows. That is the \textbf{marginal likelihood}, and
the ratio of two marginal likelihoods is the \textbf{Bayes factor}. A
flexible model spreads its predictions over many possible datasets, so
it assigns less probability to any one of them. The Bayes factor thus
implements Occam's razor automatically, and it responds to
functional-form flexibility, not just to \(k\).

Bayes factors are also used to test null hypotheses (``is sensitivity
zero?''). Unlike \(p\)-values, they can quantify evidence \textbf{for}
the null.

\paragraph{Where is this used?}

\begin{itemize}
\tightlist
\item
  \textbf{Evidence for ``no effect''.} Replication studies often ask
  whether an effect is absent. A \(p\)-value cannot support the null; a
  Bayes factor can (``the data are 8 times more likely under no
  effect'').
\item
  \textbf{Model checking in every Bayesian analysis}: posterior
  predictive checks are a standard step before any result is trusted
  (e.g.~in clinical trials analyzed with Bayesian methods and in
  cognitive modeling papers).
\item
  \textbf{Comparing theories with different flexibility}: when two
  models have the same number of parameters but one can produce many
  more data patterns, the Bayes factor is one of the few tools that
  notices.
\end{itemize}

\begin{center}\rule{0.5\linewidth}{0.5pt}\end{center}

\Needspace{11\baselineskip}

\subsection{Lecture notes}

\paragraph{In plain words}

Two separate questions this week. \textbf{(1) Is my model good enough?}
Let the fitted model \emph{imagine} new data: take plausible parameter
values from the posterior, simulate a new experiment, repeat many times.
If the real data look like a typical imagined dataset, the model is
adequate; if the real data sit outside the cloud of imagined data, the
model is missing something. This is a \textbf{posterior predictive
check}. \textbf{(2) Which of two models do the data support more?} Ask
each model, \emph{before fitting}, how probable it found the data we
actually got, averaging over all the parameter values its prior allows.
That number is the \textbf{marginal likelihood}, and the ratio for two
models is the \textbf{Bayes factor}. A model that bets on everything
(very flexible) must spread its bets thinly, so it wins less when the
data fall where a simpler model bet. That is Occam's razor, built into
probability theory. The catch: the Bayes factor depends on the prior, so
you must report how it changes with reasonable priors.

\textbf{Suggested timing (60 min):} posterior predictive checks 15 ·
marginal likelihood and Bayes factors 20 · computing BFs (grid, BIC,
Savage--Dickey) 15 · prior sensitivity and comparison of methods 10.

\subsubsection{1. Posterior predictive checks}

The \textbf{posterior predictive distribution} describes new data
\(\mathbf{y}^{\text{rep}}\) that the fitted model expects:

\[
p(\mathbf{y}^{\text{rep}} \mid \mathbf{y}) = \int p(\mathbf{y}^{\text{rep}} \mid \boldsymbol\theta)\, p(\boldsymbol\theta \mid \mathbf{y})\, d\boldsymbol\theta.
\]

\textbf{Simulation recipe:} for \(s = 1, \dots, S\), draw
\(\boldsymbol\theta^{(s)}\) from the posterior and simulate a dataset
\(\mathbf{y}^{\text{rep}(s)}\) from the model with the real design. This
propagates \emph{both} parameter uncertainty and sampling noise.

\textbf{Checks:} - \textbf{Graphical:} plot the observed data on top of
the predictive intervals (e.g., the central 95\% of \(y^{\text{rep}}\)
per condition). Systematic departures (observed points repeatedly
outside the bands, or on one side of them) show where the model fails. -
\textbf{Numerical:} choose a test statistic \(T(\mathbf{y})\) that
captures a feature you care about (e.g., recall at the longest delays,
the steepness of the early drop). The \textbf{posterior predictive
\(p\)-value} is
\(p_B = P\big(T(\mathbf{y}^{\text{rep}}) \ge T(\mathbf{y}) \mid \mathbf{y}\big)\).
Values near 0 or 1 indicate that the model cannot reproduce this
feature.

PPCs assess \textbf{absolute} adequacy, i.e., whether the model is good
enough, not which of two models is better. They are a routine step of
every Bayesian analysis (Week 13).

\paragraph{Reading the math}

\begin{itemize}
\tightlist
\item
  \(\mathbf{y}^{\text{rep}}\): ``replicated'' data, an imagined
  repetition of the experiment.
\item
  \(\int p(\mathbf{y}^{\text{rep}} \mid \boldsymbol\theta)\, p(\boldsymbol\theta \mid \mathbf{y})\, d\boldsymbol\theta\):
  average the model's predictions over all parameter values, each
  weighted by its posterior plausibility. In practice this integral is
  replaced by the two-step simulation recipe: draw \(\boldsymbol\theta\)
  from the posterior, then simulate data with it.
\item
  \(p_B = P(T(\mathbf{y}^{\text{rep}}) \ge T(\mathbf{y}) \mid \mathbf{y})\):
  the fraction of imagined datasets whose summary \(T\) is at least as
  large as the real one. If only 1\% of imagined datasets show recall as
  high as observed at 30 days (\(p_B = 0.01\)), the model cannot produce
  that feature.
\end{itemize}

\subsubsection{2. The marginal likelihood and the Bayes factor}

The \textbf{marginal likelihood} (evidence) of model \(\mathcal{M}\) is
the likelihood averaged over the prior:

\[
p(\mathbf{y} \mid \mathcal{M}) = \int p(\mathbf{y} \mid \boldsymbol\theta, \mathcal{M})\, p(\boldsymbol\theta \mid \mathcal{M})\, d\boldsymbol\theta.
\]

It is the probability the model assigned to the observed data
\emph{before seeing them}. Bayes' rule at the level of models gives:

\[
\underbrace{\frac{p(\mathcal{M}_1 \mid \mathbf{y})}{p(\mathcal{M}_2 \mid \mathbf{y})}}_{\text{posterior odds}}
= \underbrace{\frac{p(\mathbf{y} \mid \mathcal{M}_1)}{p(\mathbf{y} \mid \mathcal{M}_2)}}_{\mathrm{BF}_{12}}
\times \underbrace{\frac{p(\mathcal{M}_1)}{p(\mathcal{M}_2)}}_{\text{prior odds}}.
\]

\(\mathrm{BF}_{12}\) is the factor by which the data change the odds in
favor of \(\mathcal{M}_1\) over \(\mathcal{M}_2\), and
\(\mathrm{BF}_{21} = 1/\mathrm{BF}_{12}\).

\textbf{The automatic Occam's razor.} Both models must spread a total
probability of 1 over all possible datasets. A simple model concentrates
its predictions on a few datasets; a flexible model spreads them thin.
If the data fall where the simple model predicted, the simple model gets
the higher marginal likelihood, even though the flexible model can also
fit them well after its parameters are tuned. Unlike AIC/BIC, this
penalty depends on \emph{how} flexible a model is (the prior predictive
spread), not merely on \(k\).

\textbf{A common evidence scale} (an adaptation of Jeffreys' classic
categories, as used by Lee \& Wagenmakers):

\begin{longtable}[]{@{}ll@{}}
\caption{\centering 2. The marginal likelihood and the Bayes factor}\tabularnewline
\toprule\noalign{}
\(\mathrm{BF}_{12}\) & Evidence for \(\mathcal{M}_1\) \\
\midrule\noalign{}
\endfirsthead
\toprule\noalign{}
\(\mathrm{BF}_{12}\) & Evidence for \(\mathcal{M}_1\) \\
\midrule\noalign{}
\endhead
\bottomrule\noalign{}
\endlastfoot
1--3 & Anecdotal \\
3--10 & Moderate \\
10--30 & Strong \\
30--100 & Very strong \\
\textgreater{} 100 & Extreme \\
\end{longtable}

(Values below 1 are read the same way for \(\mathcal{M}_2\).) These
labels are aids, not thresholds. Report the number.

\paragraph{Reading the math}

\begin{itemize}
\tightlist
\item
  \(p(\mathbf{y} \mid \mathcal{M}) = \int p(\mathbf{y} \mid \boldsymbol\theta)\,p(\boldsymbol\theta)\,d\boldsymbol\theta\):
  the likelihood \textbf{averaged} over the prior, not maximized.
  Compare: MLE takes the \emph{best} parameter value; the marginal
  likelihood takes the \emph{prior-weighted average} over all of them. A
  model with many bad parameter values in its prior gets a low average.
\item
  \textbf{Posterior odds = Bayes factor × prior odds.} Odds are
  probability ratios: odds of 3 mean 3:1, i.e.~probability 0.75. If both
  models were equally plausible before (prior odds = 1) and
  \(\mathrm{BF}_{12} = 3\), then afterwards
  \(P(\mathcal{M}_1 \mid \mathbf{y}) = 3/(1+3) = 0.75\).
\item
  \textbf{Worked example in one line:} with 7 of 10 correct, chance
  predicted the data with probability 0.117 and the ``any accuracy''
  model with 0.091, so \(\mathrm{BF}_{01} = 1.29\): the data hardly
  change the odds.
\end{itemize}

\subsubsection{3. Computing Bayes factors}

\textbf{(a) Grid integration} (1--2 parameters):
\(p(\mathbf{y}\mid\mathcal{M}) \approx \sum_g p(\mathbf{y}\mid\boldsymbol\theta_g)\, p(\boldsymbol\theta_g)\, \Delta_g\),
where \(\Delta_g\) is the cell area. Do it in log space (log-sum-exp),
and check that the prior mass inside the grid is \ensuremath{\approx} 1.

\textbf{(b) The BIC approximation.} For large \(N\),
\(\ln p(\mathbf{y}\mid\mathcal{M}) \approx -\tfrac12\mathrm{BIC}\), so

\[
\mathrm{BF}_{12} \approx \exp\big(-\tfrac12(\mathrm{BIC}_1 - \mathrm{BIC}_2)\big).
\]

It is quick but crude: it ignores the actual prior (it implicitly
assumes a ``unit-information'' prior) and functional-form flexibility.

\textbf{(c) The Savage--Dickey density ratio} (nested models). If
\(\mathcal{M}_0\) fixes \(\phi = \phi_0\) in \(\mathcal{M}_1\), and the
priors on the remaining parameters match, then

\[
\mathrm{BF}_{01} = \frac{p(\phi = \phi_0 \mid \mathbf{y}, \mathcal{M}_1)}{p(\phi = \phi_0 \mid \mathcal{M}_1)}:
\]

the posterior density at the null value divided by the prior density at
the null value, both under the larger model. If the data make \(\phi_0\)
more plausible than it was a priori, the evidence favors the null. One
posterior from \(\mathcal{M}_1\) is all we need.

(Further methods, such as bridge sampling and thermodynamic integration,
compute marginal likelihoods from MCMC output for complex models. They
are beyond this course.)

\paragraph{Reading the math}

\textbf{Savage--Dickey}
\(\mathrm{BF}_{01} = p(\phi_0 \mid \mathbf{y}) / p(\phi_0)\): look at
the height of the posterior and of the prior at the null value
\(\phi_0\) (e.g.~\(\theta = 0.5\)). If the posterior is \emph{higher}
there than the prior was, the data made the null value more plausible,
so the evidence favors the null. In the worked example the prior (flat)
has height 1 at 0.5 and the posterior \(\mathrm{Beta}(8,4)\) has height
1.29, so \(\mathrm{BF}_{01} = 1.29\). A picture of two curves and one
vertical line at \(\phi_0\) is all you need to explain this.

\subsubsection{4. Prior sensitivity: estimation vs.~testing}

For \textbf{estimation}, a reasonable prior matters less and less as
data accumulate. For \textbf{Bayes factors}, the prior on a tested
parameter \emph{never} stops mattering. Making the prior under
\(\mathcal{M}_1\) wider spreads its predictions more thinly, which
lowers \(p(\mathbf{y}\mid\mathcal{M}_1)\), so an arbitrarily vague prior
makes the BF favor the null no matter what the data are (the
Jeffreys--Lindley paradox). Consequences: - Priors for model comparison
must be chosen thoughtfully, ideally reflecting realistic effect sizes.
- Always report a \textbf{sensitivity analysis}: the BF over a range of
reasonable prior widths.

\subsubsection{5. Comparison of model-evaluation methods}

\begin{longtable}[]{@{}
  >{\raggedright\arraybackslash}p{(\linewidth - 8\tabcolsep) * \real{0.2000}}
  >{\raggedright\arraybackslash}p{(\linewidth - 8\tabcolsep) * \real{0.2000}}
  >{\raggedright\arraybackslash}p{(\linewidth - 8\tabcolsep) * \real{0.2000}}
  >{\raggedright\arraybackslash}p{(\linewidth - 8\tabcolsep) * \real{0.2000}}
  >{\raggedright\arraybackslash}p{(\linewidth - 8\tabcolsep) * \real{0.2000}}@{}}
\caption{\centering 5. Comparison of model-evaluation methods}\tabularnewline
\toprule\noalign{}
\begin{minipage}[b]{\linewidth}\raggedright
Method
\end{minipage} & \begin{minipage}[b]{\linewidth}\raggedright
Question
\end{minipage} & \begin{minipage}[b]{\linewidth}\raggedright
Penalizes
\end{minipage} & \begin{minipage}[b]{\linewidth}\raggedright
Prior needed?
\end{minipage} & \begin{minipage}[b]{\linewidth}\raggedright
Week
\end{minipage} \\
\midrule\noalign{}
\endfirsthead
\toprule\noalign{}
\begin{minipage}[b]{\linewidth}\raggedright
Method
\end{minipage} & \begin{minipage}[b]{\linewidth}\raggedright
Question
\end{minipage} & \begin{minipage}[b]{\linewidth}\raggedright
Penalizes
\end{minipage} & \begin{minipage}[b]{\linewidth}\raggedright
Prior needed?
\end{minipage} & \begin{minipage}[b]{\linewidth}\raggedright
Week
\end{minipage} \\
\midrule\noalign{}
\endhead
\bottomrule\noalign{}
\endlastfoot
\(G^2\), residuals, PPC & Is the model adequate? & --- & PPC: yes & 4,
10 \\
LRT & Does the extra parameter improve fit beyond chance? (nested only)
& \(\Delta k\) via \(\chi^2\) & No & 7 \\
AIC / CV & Which model predicts new data best? & Estimated over-fitting
& No & 6--7 \\
BIC & Approximately: which model has the higher evidence? & \(k \ln N\)
& Implicit & 7 \\
Bayes factor & How much do the data change the odds between models? &
Prior predictive spread (true flexibility) & \textbf{Yes, and the result
depends on it} & 10 \\
\end{longtable}

\begin{quote}
\textbf{Advanced: Bayes factors vs.~predictive criteria.} The marginal
likelihood can be written as a product of sequential predictions,
\(p(\mathbf{y}) = p(y_1)\,p(y_2\mid y_1)\cdots p(y_N\mid y_{1:N-1})\).
It therefore scores how well the model predicts each data point from the
earlier ones, \emph{starting from the prior}. LOO-CV scores predictions
from (almost) all the other data. This is why BFs depend on the prior
and LOO hardly does, and why they can disagree on the same data.
\end{quote}

\paragraph{For an engineer}

\begin{itemize}
\tightlist
\item
  \textbf{Posterior odds = BF × prior odds} is the \textbf{MAP (Bayes)
  detector} between two hypotheses, where the Bayes factor plays the
  role of the likelihood ratio.
\item
  The difference from the \textbf{GLRT} (Week 4 and 7): with unknown
  parameters, the GLRT \emph{maximizes} the likelihood over them; the
  Bayes factor \emph{integrates} (averages) over them with a prior.
  Maximizing favors flexible hypotheses; averaging penalizes them
  automatically.
\item
  The \textbf{Jeffreys--Lindley paradox} has an engineering flavor: if
  you let the alternative hypothesis allow absurdly large signal
  amplitudes (a very wide prior), almost any observation is better
  explained by ``noise only''.
\item
  \textbf{Posterior predictive checks} are like validating a channel
  model by simulating received signals from the fitted model and
  comparing their statistics (histograms, autocorrelations, eye
  diagrams) with measured ones.
\end{itemize}

\paragraph{For the instructor}

\begin{itemize}
\tightlist
\item
  \textbf{Explain Occam's razor with a betting story.} Two forecasters
  each have 100 chips to place on possible outcomes. One spreads them
  over everything; the other puts most chips on a narrow range. If the
  outcome lands in that narrow range, the focused forecaster wins more.
  The marginal likelihood is the payout.
\item
  \textbf{Keep the two questions apart.} Write ``Is the model good
  enough? \ensuremath{\rightarrow} PPC'' and ``Which model is better? \ensuremath{\rightarrow} BF / CV / AIC'' on the
  board. Students often use a BF to claim a model is \emph{good}; it
  only says it is \emph{better than the other}.
\item
  \textbf{Typical questions.}

  \begin{itemize}
  \tightlist
  \item
    \emph{``BF or \(p\)-value?''} They answer different questions. The
    BF compares two specific hypotheses and can support the null; the
    \(p\)-value only measures surprise under the null.
  \item
    \emph{``BF = 2.5, is that significant?''} There is no threshold; it
    is anecdotal evidence. Report the number.
  \item
    \emph{``Why do BF and AIC disagree?''} They aim at different targets
    (evidence vs.~prediction) and the BF depends on the prior.
  \end{itemize}
\end{itemize}

\begin{center}\rule{0.5\linewidth}{0.5pt}\end{center}

\Needspace{11\baselineskip}

\subsection{Worked example}

\subsubsection{Part 1: is recall at chance? A Bayes factor by hand}

A participant chooses between two alternatives on \(n = 10\) trials and
is correct on \(y = 7\). Is performance at chance?

\begin{itemize}
\tightlist
\item
  \(\mathcal{M}_0\): \(\theta = 0.5\) (no free parameters).
\item
  \(\mathcal{M}_1\): \(\theta \sim \mathrm{Beta}(1, 1)\) (any accuracy
  equally plausible a priori).
\end{itemize}

\textbf{Step 1.}
\(p(y\mid\mathcal{M}_0) = \binom{10}{7} 0.5^{10} = 120/1024 = 0.1172\).

\textbf{Step 2.} Under a uniform prior, every count \(0, \dots, n\) is
equally likely a priori, so
\(p(y\mid\mathcal{M}_1) = \int_0^1 \binom{10}{7}\theta^7(1-\theta)^3\, d\theta = \frac{1}{n+1} = \frac{1}{11} = 0.0909\).

\textbf{Step 3.} \(\mathrm{BF}_{01} = 0.1172 / 0.0909 = 1.29\). The data
slightly \emph{favor chance}, which is anecdotal evidence either way.

\textbf{Step 4 (Savage--Dickey check).} The posterior under
\(\mathcal{M}_1\) is \(\mathrm{Beta}(8, 4)\), whose density at 0.5 is
1.289. The prior density at 0.5 is 1. The ratio is 1.29, identical to
Step 3.

\textbf{More data:}

\begin{longtable}[]{@{}
  >{\raggedright\arraybackslash}p{(\linewidth - 6\tabcolsep) * \real{0.2500}}
  >{\raggedright\arraybackslash}p{(\linewidth - 6\tabcolsep) * \real{0.2500}}
  >{\raggedright\arraybackslash}p{(\linewidth - 6\tabcolsep) * \real{0.2500}}
  >{\raggedright\arraybackslash}p{(\linewidth - 6\tabcolsep) * \real{0.2500}}@{}}
\caption{\centering Part 1: is recall at chance? A Bayes factor by hand}\tabularnewline
\toprule\noalign{}
\begin{minipage}[b]{\linewidth}\raggedright
Data
\end{minipage} & \begin{minipage}[b]{\linewidth}\raggedright
\(p\) (two-sided binomial test)
\end{minipage} & \begin{minipage}[b]{\linewidth}\raggedright
\(\mathrm{BF}_{01}\)
\end{minipage} & \begin{minipage}[b]{\linewidth}\raggedright
Reading
\end{minipage} \\
\midrule\noalign{}
\endfirsthead
\toprule\noalign{}
\begin{minipage}[b]{\linewidth}\raggedright
Data
\end{minipage} & \begin{minipage}[b]{\linewidth}\raggedright
\(p\) (two-sided binomial test)
\end{minipage} & \begin{minipage}[b]{\linewidth}\raggedright
\(\mathrm{BF}_{01}\)
\end{minipage} & \begin{minipage}[b]{\linewidth}\raggedright
Reading
\end{minipage} \\
\midrule\noalign{}
\endhead
\bottomrule\noalign{}
\endlastfoot
7/10 & 0.34 & 1.29 & Anecdotal, either way \\
60/100 & \ensuremath{\approx} 0.06 & 1.10 & Still anecdotal: the data barely distinguish
chance from ``some accuracy'' \\
70/100 & \textless{} 0.001 & 0.0023 (\(\mathrm{BF}_{10} \approx 430\)) &
Extreme evidence against chance \\
\end{longtable}

\textbf{Interpretation.} With 60/100, a \(p\)-value near .05 tempts
researchers to claim an effect, yet the Bayes factor says the data are
about equally likely under both hypotheses. The reason: the alternative
with a uniform prior ``predicted'' that accuracy could be anywhere from
0 to 1. An observed accuracy of 0.6 is only modestly better predicted by
it than by exactly 0.5. A narrower, more realistic prior under
\(\mathcal{M}_1\) (e.g., accuracies between 0.5 and 0.8) would give a
different BF, which is why prior sensitivity must be reported.

\subsubsection{Part 2: posterior model probabilities}

For 7/10 with prior odds of 1: posterior odds \(= 1.29\), so
\(P(\mathcal{M}_0 \mid y) = 1.29 / 2.29 = 0.56\). If a reviewer believes
a priori that chance performance is 4 times less likely than
above-chance performance (prior odds 1:4), then posterior odds
\(= 1.29 \times 0.25 = 0.32\), so \(P(\mathcal{M}_0 \mid y) = 0.24\).
The BF is the same for everyone; the posterior model probabilities
depend on the prior odds.

\begin{center}\rule{0.5\linewidth}{0.5pt}\end{center}

\Needspace{11\baselineskip}

\subsection{Lab (2 hours)}

\subsubsection{Goals}

\begin{enumerate}
\def\labelenumi{\arabic{enumi}.}
\tightlist
\item
  Run graphical and numerical posterior predictive checks for the
  exponential and power forgetting models.
\item
  Compute marginal likelihoods and the Bayes factor for power
  vs.~exponential by grid integration; compare with BIC.
\item
  Test \(d' = 0\) for a weak-evidence SDT participant with the
  Savage--Dickey ratio, and verify it by direct computation.
\item
  Map the sensitivity of the Bayes factor to the prior, and contrast it
  with the posterior estimate.
\end{enumerate}

\textbf{Timing:} Part A 35 min · Part B 25 min · Part C 30 min · Part D
20 min · Part E 10 min.

\subsubsection{Part A: posterior predictive checks (35 min)}

\textbf{Interpret:} look at the plots first. For the exponential model,
the observed recall at the shortest delay (42/50) lies \emph{above} the
95\% band and recall at 7 days (13/50) lies \emph{below} it. The
exponential cannot drop steeply early and then level off. The power
model's band contains every observed point.

The exponential's misfit is in the \emph{shape} (early vs.~middle
delays), which this statistic does not look at. \textbf{A PPC can only
detect the failures its test statistic is sensitive to.} Choose
statistics that target the features that distinguish the theories; in
Exercise 2 a ``first drop'' statistic exposes the exponential model
(\(p_B \approx 0.02\)).

\subsubsection{Part B: Bayes factor by grid integration (25 min)}

\textbf{Interpret:} \(\mathrm{BF} \approx 1.7 \times 10^5\) is extreme
evidence for the power model over the exponential. BIC gives the same
verdict but a different number (\(2 \times 10^4\)). Why do the grid BF
and the BIC approximation differ? (Hint: which prior does each use, and
are the two models equally flexible under our priors?)

\subsubsection{Part C: Savage--Dickey: is this participant's sensitivity
zero? (30 min)}

A participant in a hard detection task: 12 hits out of 20 signal trials,
8 false alarms out of 20 noise trials. \(\mathcal{M}_1\):
\(d' \sim \mathcal{N}(0, 1)\), \(c \sim \mathcal{N}(0, 1)\).
\(\mathcal{M}_0\): \(d' = 0\), \(c \sim \mathcal{N}(0, 1)\).

The two computations agree, which is a useful check that both are
implemented correctly.

\subsubsection{Part D: prior sensitivity (20 min)}

\textbf{What to notice:} as the prior on \(d'\) widens, the posterior
mean of \(d'\) quickly stabilizes, but \(\mathrm{BF}_{01}\) keeps
growing in favor of the null. A very wide prior ``predicts'' huge
sensitivities that did not occur, so the alternative is penalized.
\textbf{For estimation the prior washes out; for testing it does not.}

\subsubsection{Part E: interpret (10 min)}

Write 4--5 sentences: does this participant detect the signal? Report
the BF with your chosen prior and justify the prior, mention the
sensitivity analysis, and state what additional data would make the
answer clearer.

\subsubsection{Exercises (increasing difficulty)}

\(\int_0^1 \binom{10}{7}\theta^7(1-\theta)^3 d\theta = 1/11\)

\textbf{Exercise 2 \ensuremath{\star}.} In Part A, use a different test statistic: the
drop in recall from the first to the second interval. Which model, if
any, fails this check? (Expected: the exponential, \(p_B \approx 0.02\);
the power model passes.)

\textbf{Exercise 3 \ensuremath{\star}\ensuremath{\star} (prior sensitivity of the forgetting BF).} Repeat
Part B with the rate priors widened (scale × 3) and narrowed (scale ÷
3). Does the conclusion ``power \ensuremath{\gg} exponential'' survive? Why is it so
robust here, unlike in Part D?

\textbf{Exercise 4 \ensuremath{\star}\ensuremath{\star} (an informed alternative).} In Part C, replace the
prior on \(d'\) by a half-normal \(d' \sim \mathcal{N}^+(0, 1)\) (only
positive values; remember to renormalize). Compute the Savage--Dickey
BF. Why does a directional prior change the result?

\textbf{Exercise 5 \ensuremath{\star}\ensuremath{\star}\ensuremath{\star} (Advanced: BF for the repetition study).} For the
Week 8 case study (\(\mathcal{M}_0\) vs.~\(\mathcal{M}_1\)), choose
priors for \(\delta_0\), \(\delta_1\) and \(\kappa\), compute
\(\mathrm{BF}_{10}\) with the Savage--Dickey ratio at \(\delta_1 = 0\)
using a 3-D grid (keep it coarse), and compare it with the BIC
approximation.

\subsubsection{Lab check-off}

Show: (1) the two PPC plots with \(p_B\) values, (2) the grid BF for
power vs.~exponential and its BIC approximation, and (3) the
prior-sensitivity table from Part D with one sentence of interpretation.

\begin{center}\rule{0.5\linewidth}{0.5pt}\end{center}

\subsection{Connection to AI \& LLMs}

\textbf{Do LLMs have a ``theory of mind''?} Some studies report that
LLMs pass false-belief tasks, classic tests of whether one understands
that others can hold false beliefs. Others show that small changes to
the stories make performance collapse {[}verify: Kosinski, 2023/2024;
Ullman, 2023{]}. Many such tests have two answer options, so chance is
50\%, and the number of items is small. A Bayes factor answers the right
first question: \emph{do the data support above-chance performance at
all, compared with chance?}

The same 65\% is no evidence at all with 40 items, and overwhelming
evidence with 200 items. Above-chance performance is also only the first
step: the posterior predictive check of this week asks \emph{how} the
model succeeds and fails. For example, does it fail on exactly the
altered stories, where a human would not?

\(\theta > 0.5\) How does \(\mathrm{BF}_{10}\) for 26/40 change? - \ensuremath{\star}\ensuremath{\star}
Simulate a ``shortcut'' model that answers correctly on 90\% of standard
stories and 30\% of altered stories (half of each). Its overall accuracy
is 60\%. Design a posterior predictive check that reveals the shortcut.
- \ensuremath{\star}\ensuremath{\star} Discussion: is ``above chance on false-belief tasks'' a claim at
Marr's computational or algorithmic level?

\begin{center}\rule{0.5\linewidth}{0.5pt}\end{center}

\subsection{Common mistakes \& misconceptions}

\begin{enumerate}
\def\labelenumi{\arabic{enumi}.}
\tightlist
\item
  \textbf{``A narrow posterior means the model is right.''} Precision is
  not adequacy. Always run a PPC.
\item
  \textbf{Using the posterior predictive check to choose between
  models.} PPCs check adequacy. Use BFs, CV or information criteria to
  compare.
\item
  \textbf{``The Bayes factor is the posterior probability of the
  model.''} It is the \emph{change} in odds. Posterior probabilities
  also need prior odds.
\item
  \textbf{Using very vague priors in Bayes-factor tests.} They bias the
  BF toward the null without limit (Part D). Use realistic priors and
  report the sensitivity.
\item
  \textbf{Savage--Dickey with mismatched priors.} The nuisance-parameter
  priors must be the same in both models, and the prior density must be
  evaluated \emph{at the null value} in the larger model.
\item
  \textbf{Computing marginal likelihoods on a grid that cuts off prior
  mass.} Check that the prior mass inside the grid is close to 1.
\item
  \textbf{Treating ``BF = 2.9'' vs.~``BF = 3.1'' as qualitatively
  different.} The evidence scale is continuous; labels are aids, not
  decision thresholds.
\end{enumerate}

\begin{center}\rule{0.5\linewidth}{0.5pt}\end{center}

\subsection{Check your understanding}

\textbf{1.} How is a posterior predictive distribution simulated, and
what does it combine? \textbf{Answer} Draw parameters from the
posterior, then simulate data from the model with those parameters (and
the real design). It combines uncertainty about the parameters with the
sampling noise of new data.

\textbf{2.} Define the marginal likelihood in words. \textbf{Answer} The
probability (density) of the observed data under the model, averaged
over all parameter values weighted by the prior: how well the model
predicted these data before it was fitted.

\textbf{3.} Why does the marginal likelihood penalize flexible models
even without counting parameters? \textbf{Answer} Probability over
possible datasets must sum to 1. A flexible model spreads its prior
predictions over many datasets and so assigns each one, including the
observed one, less probability, unless the data fall where only the
flexible model predicts them.

\textbf{4.} \(\mathrm{BF}_{10} = 6\) and the prior odds are 1:3 (against
\(\mathcal{M}_1\)). What are the posterior odds and
\(P(\mathcal{M}_1\mid\mathbf{y})\)? \textbf{Answer} Posterior odds
\(= 6 \times 1/3 = 2\), so
\(P(\mathcal{M}_1\mid\mathbf{y}) = 2/3 \approx 0.67\).

\textbf{5.} State the Savage--Dickey ratio and the conditions for using
it. \textbf{Answer}
\(\mathrm{BF}_{01} = p(\phi_0\mid\mathbf{y},\mathcal{M}_1)/p(\phi_0\mid\mathcal{M}_1)\):
the posterior over the prior density at the null value under the larger
model. It requires nested models (the null fixes \(\phi = \phi_0\)) and
identical priors on the other parameters.

\textbf{6.} Why does widening the prior on the tested parameter increase
\(\mathrm{BF}_{01}\), even with fixed data? \textbf{Answer} A wider
prior assigns more probability to large effects and less to the small
effect sizes compatible with the data, so the alternative's marginal
likelihood decreases while the null's does not change.

\textbf{7.} BIC gives \(\mathrm{BF}_{12} \approx 20\), while the grid
computation with your priors gives 3. Which do you report and why?
\textbf{Answer} The grid BF, because it uses the actual priors and
accounts for functional-form flexibility; BIC is a large-sample
approximation with an implicit prior. Mention the discrepancy and,
ideally, a prior sensitivity analysis.

\begin{center}\rule{0.5\linewidth}{0.5pt}\end{center}

\Needspace{11\baselineskip}

\subsection{Assignment}

\subsubsection{A4: A Bayesian analysis by grid (1.0 point of 20;
released Week 10, due at the start of the Week 12 lecture)}

\(n = 20\)

\textbf{Tasks.}

\begin{enumerate}
\def\labelenumi{\arabic{enumi}.}
\tightlist
\item
  \textbf{Priors (1.5 pts).} Choose and justify priors for
  \((A, \lambda)\) of the exponential model and \((A, d)\) of the power
  model. Show prior predictive simulations of forgetting curves for both
  models and revise the priors if needed.
\item
  \textbf{Posterior (2 pts).} Compute the grid posteriors for both
  models. Report posterior means, SDs and 95\% HDIs, and the posterior
  correlation of the two parameters. Plot the joint posteriors.
\item
  \textbf{Posterior predictive checks (2 pts).} Graphical PPC for both
  models, plus a numerical PPC with a test statistic you choose and
  justify.
\item
  \textbf{Bayes factor (2 pts).} Compute
  \(\mathrm{BF}_{\text{pow,exp}}\) by grid integration (check the prior
  mass), compare it with the BIC approximation, and repeat with two
  alternative prior settings (sensitivity analysis).
\item
  \textbf{Comparison with A2 (1 pt).} Compare the posterior summaries
  with your A2 MLEs and their uncertainty.
\item
  \textbf{Interpretation (1.5 pts, max. 300 words).} What do you
  conclude about forgetting in these data, and how robust is the
  conclusion to the prior?
\end{enumerate}

\textbf{Rubric:} points as above; within each task: correctness 60\%,
figures and tables 20\%, interpretation 20\%. Late submissions lose 10\%
per day (max. 3 days).

\begin{center}\rule{0.5\linewidth}{0.5pt}\end{center}

\subsection{Further reading}

\begin{itemize}
\tightlist
\item
  \textbf{Lee \& Wagenmakers (2014), chapter 7.} Bayesian model
  comparison: marginal likelihoods, Bayes factors and the Savage--Dickey
  method, with worked examples. \emph{(Main reading.)}
\item
  \textbf{Farrell \& Lewandowsky (2018), chapter 11.} Bayesian model
  comparison using Bayes factors, in the context of cognitive models.
\item
  \textbf{Kruschke (2014), chapter 10.} Model comparison and
  hierarchical modeling, a gentle introduction to Bayes factors and
  their prior sensitivity.
\item
  \textbf{Martin, Kumar \& Lao (2021), chapter 2.} Exploratory analysis
  of Bayesian models, including posterior predictive checks.
\end{itemize}

\emph{Preview for Week 11:} our grids worked for two parameters. For
five, ten or a hundred parameters we need a smarter way to explore the
posterior: Markov chain Monte Carlo. You will build your first sampler
from scratch.
